\InputIfFileExists{../global.src}\relax\relax
\iffull
\title[Allgemeine Informationen]{Ja sind sie denn schon wieder da?}
\subtitle{Organisatorisches Tutorium}
\date{KW 17}
\addbibresource{references.bib}
\fi
% \SetTutoriumNumber{Org}
\setfootmarker{EidI~\faAngleRight\ Org. Tutorium}
\titlenumber{Org}
\setauthormailsubject{Organisatorisches\%20Tutorium}

% \documentclass{article}
% \usepackage{fetamont}
% \begin{document}
% \ffmwfamily\fontseries{eb}\selectfont Gauß
% \end{document}

\newsavebox\pinguA \newsavebox\pinguB
\iffull

\begin{document}
\begin{frame}[plain,c]
\begin{tikzpicture}[remember picture, overlay]
    \node at(0,0) {TODO: Übersicht};
\end{tikzpicture}
\end{frame}

\titleframe
\fi

\SetNextSectionText{Schön ist das, was ohne Begriff allgemein gefällt.\\--- Immanuel Kant~\cite[p.~61]{kant1924kritik}}
\section{Allgemeines}


\savebox\pinguA{\tikz\pingu[wings shock,eyes shock,heart=gray!25!white];}
\savebox\pinguB{\tikz\pingu[wings wave,eyes wink,mask,heart=gray!25!white,halo];}
\subsection{Hygiene}
\begin{frame}[c]{Hygienekonzept}
\begin{itemize}[<+(1)->]
    \itemsep=\fill
    \item Bisher keine besonderen Regelungen.
    \item Dennoch gilt die Empfehlung:
\end{itemize}
\vfill
\begin{center}
    \scalebox{.8}{\begin{tikzpicture}
        \contourlength{2pt}
        \onslide<4->{\node[fill=btdm@border@down,minimum size=4cm,rounded corners=5pt] (wrongy) at (0,0) {\copy\pinguA};
        \node[gray,scale=1.35] at(wrongy.south east) {\huge\contour{btdm@background}{\faTimes}};}
        \onslide<5->{\node[fill=btdm@border@down,minimum size=4cm,rounded corners=5pt,right=1.5cm] (righty) at (wrongy.east) {\raisebox{1.4pt}{\copy\pinguB}};
        \node[gray,xshift=-1mm,yshift=1mm,scale=1.35,btdm@primary] at(righty.south east) {\huge\contour{btdm@background}{\faCheck}};}
    \end{tikzpicture}}
\end{center}
\vspace*{-.75\baselineskip}
\end{frame}

\begin{frame}[t]{Folien und Tutorien}
    \begin{itemize}[<+(1)->]
        \itemsep=18pt
        \item Folien von Florian Sihler (\link{mailto:florian.sihler@uni-ulm.de}{florian.sihler@uni-ulm.de})
        \item Mitschriften, Zusammenfassungen und vieles mehr im Cloudstore: \begin{itemize}[<1->]
            \item \url{https://cloudstore.uni-ulm.de/apps/circles/}
            \item Kreis: \say{Mitschriebe Informatik} (ganz eingeben)
        \end{itemize}
        \item Bei erbrachter Leistung: 3 ECTS (LP)
    \end{itemize}
\end{frame}


\savebox\pinguA{\tikz\pingu[left wing hug,eyes wink];}
\savebox\pinguB{\tikz\pingu[right wing hug,eyes shiny];}
\subsection{Übungsblätter}
\begin{frame}[t]{Übungsblätter\hfill I}
\begin{itemize}[<+(1)->]
    % \itemsep=\fill % yielded wrong results after closing of itemize
    \item Mindestens \qty{50}\percent\ der Punkte pro Blatt\begin{itemize}
        \item Maximal 2 der 13 regulären Blätter dürfen nicht bestanden werden
        \item Aktuell keine Zusatzaufgaben geplant\vfill
    \end{itemize}
    \item Tutoriumsteilnahme verpflichtend für Punkte\vfill
    \item Abgabe in zweier Gruppen \only<handout:0|5>{\smash{\raisebox{-2pt}{\resizebox*!{\baselineskip}{\usebox\pinguA\!\usebox\pinguB}}}}\only<6->{\smash{\raisebox{-2pt}{\resizebox*!{\baselineskip}{\usebox\pinguA\hskip8em\usebox\pinguB}}}}\vfill
    \item<7-> Erlaubte Formate: {\begin{itemize}
        \item<8-> Bei einer Datei: PDF, Plaintext (\T{.txt}), Java-Quellcode
        \item<9-> Bei mehreren: Archiv (\textsc{zip},~\ldots) verwenden% TODO: NOTE: wird vermutlich noch Freigeschaltet
    \end{itemize}}\vfill
\end{itemize}
\end{frame}

\begin{frame}[t]{Übungsblätter\hfill II}
    \begin{itemize}[<+(1)->]
        \itemsep12pt
        \item Programme \emph{müssen} lauffähig sein (inklusive \bjava{main})\pause
        \infoblock{Oder entsprechend kommentiert, wo es (wieso) Probleme gibt.}
        \item \emph{Jeder} muss die Abgabe vorstellen können.\pause\infoblock{Unabhängig davon, wer die Aufgabe im Team gemacht hat.}
        \item Im Falle einer Krankheit:\pause\\\emph{Nachricht an den Tutor/die Tutorin mit zugehöriger Entschuldigung!}
    \end{itemize}
    \vfill
    \begin{center}
        \onslide<8->{\bfseries Lest die \link{https://www.uni-ulm.de/einrichtungen/zuv/dez1/recht-und-organisation/satzungen-und-ordnungen/studium-promotion-habilitation/studien-und-pruefungsordungen-bachelor-master-staatsexamen/}{\textbf{FSPO}}!\\}
        \onslide<9->{\bfseries Verwendet die Uni-Mail!}
    \end{center}
\end{frame}

\iffull
\begin{frame}[t]{Präsenzaufgaben}
    \begin{itemize}[<+(1)->]
        \itemsep10pt
        \item Stift und Papier
        \item Jede Woche, unbepunktet
        \item Darf aber gern abgegeben werden!
    \end{itemize}
\end{frame}
\fi

\begin{frame}{Klausur}
    \begin{itemize}
        \itemsep16pt
        \item Alle Informationen werden rechtzeitig bekanntgegeben
        \item Früh genug anmelden!
        \item Termine: \begin{itemize}
            \item Erstklausur am 10. August 2022
            \item Zweitklausur am 7. Oktober 2022
        \end{itemize}
    \end{itemize}
\end{frame}

\iffull
% guard against problems with wget clearing the file even if there is no connection
\immediate\write18{ping -c1 8.8.8.8 && wget https://media.githubusercontent.com/media/EagleoutIce/Episode-Heaps/gh-pages/preview-01.png -O logoHeaps-\jobname.png}
\immediate\write18{ping -c1 8.8.8.8 && wget https://media.githubusercontent.com/media/EagleoutIce/Episode-Recursion/gh-pages/preview-01.png -O logoRecursion-\jobname.png}
\immediate\write18{ping -c1 8.8.8.8 && wget https://media.githubusercontent.com/media/EagleoutIce/Episode-Traversierung/gh-pages/preview-01.png -O logoTraversal-\jobname.png}
% and more :D

{%
\IfBtdmDarkmode{
\UpdateBtdmColor*{btdm@primary}{\colorlet{btdm@block}{btdm@block!20!btdm@border@up}}
}{
\UpdateBtdmColor{btdm@primary!68!white}
}
\subsection{Zusätzliches}
\begin{frame}[t]{Zusätzliches}
    \begin{itemize}[<+(1)->]
        \itemsep9pt
        \item Es \textsl{empfiehlt} sich die Abgaben in \LaTeX\ ($2_{\textstyle\varepsilon}$) zu machen
        \item Dafür wurden erklärende Bonusblöcke erstellt. Darunter:\pause{}
        \vspace*{-.75\baselineskip}\begin{multicols}{3}
            \begin{enumerate}[<1->]
                \item \LaTeX{} (2-3 Tutorien)
                \item Richtig Kommentieren
                \item Linux (\& Kommandozeile)
                \item Git --- Versionsverwaltung
                \item (Unit-)Testen
                \item \ldots\par
            \end{enumerate}
        \end{multicols}\vspace*{-.75\baselineskip}
        \item Mit dem letzten Semester gibt es auch tolle Episoden. Darunter:\\\smallskip
% hella hail the copy and paste
\resizebox{.85\linewidth}!{\begin{tikzpicture}[align-base]
    \onslide<2->{\draw[thick,darkgray,rounded corners=2.5pt,path picture={\node at(path picture bounding box.center) {\href{https://media.githubusercontent.com/media/EagleoutIce/Episode-Traversierung/gh-pages/noanim_traversal.pdf}{\includegraphics[width=8.5cm,height=4.788cm,keepaspectratio]{logoTraversal-\jobname.png}}};}] (0,0) rectangle (8.5cm,4.788cm);}
\end{tikzpicture}\quad\begin{tikzpicture}[align-base]
    \onslide<2->{\draw[thick,darkgray,rounded corners=2.5pt,path picture={\node at(path picture bounding box.center) {\href{https://media.githubusercontent.com/media/EagleoutIce/Episode-Recursion/gh-pages/noanim_rekursion.pdf}{\includegraphics[width=8.5cm,height=4.788cm,keepaspectratio]{logoRecursion-\jobname.png}}};}] (0,0) rectangle (8.5cm,4.788cm);}
\end{tikzpicture}\quad\begin{tikzpicture}[align-base]
    \onslide<2->{\draw[thick,darkgray,rounded corners=2.5pt,path picture={\node at(path picture bounding box.center) {\href{https://media.githubusercontent.com/media/EagleoutIce/Episode-Heaps/gh-pages/noanim_heap.pdf}{\includegraphics[width=8.5cm,height=4.788cm,keepaspectratio]{logoHeaps-\jobname.png}}};}] (0,0) rectangle (8.5cm,4.788cm);}
\end{tikzpicture}\qquad\textbullet~\textbullet~\textbullet}
        \item Die (Lösungs-)Folien werden wochenweise veröffentlicht.
    \end{itemize}
\end{frame}
}
\fi

\SetNextSectionText[.7\linewidth]{As long as there were no machines, programming was no problem at all; when we had a few weak computers, programming became a mild problem, and now we have gigantic computers, programming had become an equally gigantic problem. In this sense the electronic industry has not solved a single problem, it has only created them.\\--- Edsger W. Dijkstra~\cite[p.~861]{dijkstra1972humble}}
\section{Java \& Editoren}
\def\os#1{\;\textsuperscript{\tiny\color{btdm@text}\def\x##1{\textcolor{gray!80!lightgray}{##1}}#1}}
\subsection{Was es für Editoren gibt}
\begin{frame}[t]{How to Editor}
    \begin{itemize}[<+(1)->]
        \itemsep16pt
        \item Text-Editor mit Syntaxfreude
        \item Keine IDE!
        \item Klassisch: \begin{itemize}
            \item \link{https://notepad-plus-plus.org/}{Nodepad++}\os{\faWindows, \x{\faApple}, \x{\faLinux}}
            \item \link{https://code.visualstudio.com/Download}{VS-Code}\os{\faLinux, \faApple, \faWindows} (ohne Extensions)
            \item \link{https://atom.io/}{\textsc{atom}}\os{\faLinux, \faApple, \faWindows}
        \end{itemize}
        \item Advanced: \begin{itemize}
            \item \link{https://www.nano-editor.org/}{nano}\os{\faLinux, \faApple, \x{\faWindows}}
            \item \link{https://www.vim.org/}{vim}\os{\faLinux, \faApple, \x{\faWindows}}
        \end{itemize}
    \end{itemize}
\end{frame}

\subsection{Wie Java geht}
\begin{frame}[t]{How to Java}
    \begin{itemize}[<+(1)->]
        \itemsep5.5pt
        \item Linux: \begin{itemize}
            \item Wenn \T{apt}-basiert: \bbash{sudo apt install default-jdk} (\T{openjdk-14-jdk},~\ldots)
            \item Auch sonst über den jeweiligen Paketmanager
            \item Oder per \link{https://www.oracle.com/java/technologies/javase-downloads.html}{download}.
        \end{itemize}
        \item Windows: \begin{itemize}
            \item Download \& Installation der \link{https://www.oracle.com/java/technologies/javase-downloads.html}{gewünschten Version}.
            \item \T{PATH}-Variable anpassen! Dafür gibt es etliche \link{https://explainjava.com/set-java-path-and-java-home-windows/}{Anleitungen}:
                \begin{enumerate}
                    \item \say{Systemumgebungsvariablen bearbeiten} beziehungsweise \say{Edit the system environment variables}
                    \item \say{bin}-Ordner der JDK zum \T{PATH} hinzufügen. \info{Zum Beispiel: \say{C:\textbackslash Program Files\textbackslash Java\textbackslash jdk-15\textbackslash bin}}
                \end{enumerate}
        \end{itemize}
        \item MacOS: \begin{itemize}
            \item Download \& Installation der \link{https://www.oracle.com/java/technologies/javase-downloads.html}{gewünschten Version}.
            \item Sonst unterstützt auch \link{https://brew.sh/}{Homebrew} ein paar Versionen.
        \end{itemize}
        \item Es existieren \link{https://jdk.java.net/16/}{open source} Implementationen.
    \end{itemize}
\end{frame}

\SetNextSectionText{Gefühl von Grenze darf nicht heißen: hier bist du zu Ende, sondern: hier hast du noch zu wachsen.\\--- Emil Gött~\cite{gott1984zettelspruche}}
\section{Abschließendes}
{\SummaryFrame
\begin{frame}[t]{Zusammenfassend}
\pause \printBibCommand
\vfill\vfill % double fill for more fraction
\begin{itemize}[<+(1)->]
    \itemsep16pt
    \item Wöchentliche Übungsblätter in Zweiergruppen \begin{itemize}
        \item {\qty{50}\percent} der Punkte und Anwesenheit im Tutorium um Blatt zu Bestehen.
        \item Mindestens 11 der 13 Blätter müssen bestanden werden (exklusive Blatt 0).
        \item Beide müssen die Abgabe vorstellen können.
        \item Programme müssen lauffähig sein.
    \end{itemize}
    \item Bei Fragen darf sich jederzeit \emph{gerne} gemeldet werden
\end{itemize}
\end{frame}
}

\outro{\begin{tikzpicture}% remember picture does not work :C
\IfBtdmDarkmode{}{\pgfinterruptboundingbox
\only<3->{\fill[btdm@primary] (0,0) circle [radius=2*\paperwidth];}
\endpgfinterruptboundingbox}
\end{tikzpicture}%
\null\vfill
% use fill and stuff to get it done
\centering
\only<2|handout:0>{\begin{tikzpicture}[scale=1.5]
    \pingu[wings wave,name=saphira,eyes wink,pants=cprimary,monocle right]
    \pgfinterruptboundingbox
    \onslide<2->{\path[postaction={decorate},decoration={text along path, text={|\huge|Motivation!},text align={fit to path}}] (saphira-bill) ++ (90+55:45pt) arc (90+55:90-55:45pt);}
    \endpgfinterruptboundingbox
\end{tikzpicture}}%
\only<3->{%
\IfBtdmDarkmode{}{\BtdmMakeNextFullBottom}%inverts bottom
\begin{tikzpicture}[scale=1.5]
    \pingu[wings raise,monocle right,name=saphira,eyes wink,devil horns,left eye devil,right eye angry,bill=angry,glow thick=white,heart=pingu@purple!60!black]
    \pgfinterruptboundingbox
    \path[postaction={decorate},decoration={text along path, text={|\huge\bfseries\IfBtdmDarkmode{\color{btdm@primary}}{\color{white}}|Motivation!},text align={fit to path}}] (saphira-bill) ++ (90+55:45pt) arc (90+55:90-55:45pt);
    \endpgfinterruptboundingbox
\end{tikzpicture}%
}%
}
% {\setbeamercolor{background canvas}{bg=cprimary!15!black}
% \begin{frame}[plain,c]
% \null\vfill
% \end{frame}}
\iffull\end{document}\fi